% sections/06_scope_limits.tex
% JPHQ-04 — Scope limits and non-claims (OC/ETS-internal)

\subsection*{Scope Limits}

The No-Go theorems stated in this preprint are \emph{conditional} results.
They assert impossibility of enterprise transformation only under explicitly
stated structural constraints within the \OC/\kw{ETS.K\_EA.v1.1} formal system.
No claim of universal, unconditional, or context-independent impossibility is
made.

The scope of the results is restricted as follows:
\begin{enumerate}
  \item All arguments and derivations are internal to the assumed OC/ETS
  primitives, operators, and admissibility criteria.
  \item Admissibility is evaluated exclusively via:
  (i) existence of states in $\Om(\KEA)$,
  (ii) existence of admissible cycles in $\Cyc(\KEA)$, and
  (iii) non-vanishing continuity $\kcont(\KEA)>0$.
  \item Phase labels $\PHOK$, $\PHIN$, $\PHCO$, and $\PHST$ are treated strictly
  as diagnostic classifications of structural admissibility and continuity,
  without evaluative, normative, teleological, or managerial interpretation.
\end{enumerate}

Accordingly, the No-Go theorems delineate regions of \emph{formal non-existence}:
classes of structural configurations for which transformation cannot occur
without violating admissibility, collapsing the admissible state space, or
entering terminal phases as defined by the OC/ETS phase logic.

\subsection*{Non-Claims}

This artifact makes no prescriptive, managerial, or normative claims. In
particular, it does \emph{not} provide:
\begin{enumerate}
  \item recommendations, interventions, or governance frameworks;
  \item optimization objectives, control strategies, or steering mechanisms;
  \item performance measures, KPIs, benchmarks, or empirically calibrated models;
  \item narratives of success, recovery, best practice, or improvement;
  \item teleological statements concerning what enterprises \emph{should} do.
\end{enumerate}

The sole purpose of this section is to prevent category errors. The No-Go
theorems are statements about formal impossibility under specified structural
conditions within \OC/\kw{ETS}. They must not be interpreted as statements about
desirability, strategic choice, feasibility of intervention, or actionability.
